\section*{Capítulo \ref{ch:g9:atoms-to-galaxies} --- Dos átomos às galáxias: as escalas do Universo}

\begin{solution}{exo:g9:atoms-to-galaxies:1}
\omterm{def:g9:atoms-to-galaxies:atom}{Núcleo} ($10^{-15}$ \unit{m}), \omterm{def:g9:atoms-to-galaxies:atom}{átomo} ($10^{-10}$), \omterm{def:g7:states-of-matter:molecule}{molécula} (em torno
de $10^{-9}$), formiga ($10^{-2}$, ou $10^{-3}$ para uma modesta).
\end{solution}

\begin{solution}{exo:g9:atoms-to-galaxies:2}
Um \omterm{def:g9:atoms-to-galaxies:atom}{núcleo} minúsculo, pesado e de carga positiva no centro; uma nuvem
de \omterm{def:g9:atoms-to-galaxies:atom}{elétrons} negativos e quase sem \omterm{def:g4:weight-and-mass:weight}{peso} bem longe dele. Os tamanhos:
\omterm{def:g9:atoms-to-galaxies:atom}{átomo} perto de $10^{-10}$ \unit{m}, \omterm{def:g9:atoms-to-galaxies:atom}{núcleo} perto de $10^{-15}$ --- o
abismo de cem mil vezes que faz a ervilha num estádio.
\end{solution}

\begin{solution}{exo:g9:atoms-to-galaxies:3}
Os \omterm{def:g9:atoms-to-galaxies:atom}{elétrons} --- negativos, marchando do $-$ para o $+$: ao contrário
da convenção. As leis sobrevivem porque nunca dependeram de quem
marcha: as correntes, as somas nos \omterm{def:g5:series-parallel-circuits:parallel}{nós}, os retratos de Ohm e as
contas de \omterm{def:g1:light-and-shadows:source}{luz} se leem de modo idêntico com a seta mantida por acordo
mundial.
\end{solution}

\begin{solution}{exo:g9:atoms-to-galaxies:4}
A potência de dez mais próxima. Altura: $10^{0}$ \unit{m}; sala de
aula: $10^{1}$; diâmetro da Terra: $10^{7}$.
\end{solution}

\begin{solution}{exo:g9:atoms-to-galaxies:5}
Do \omterm{def:g9:atoms-to-galaxies:atom}{átomo} ao \omterm{def:g9:atoms-to-galaxies:atom}{núcleo}: cinco degraus. De você à Terra: sete. Da
distância do Sol ($10^{11}$) à da \omterm{def:g4:solar-system:star}{estrela} mais próxima
($4 \times 10^{16}$): cinco a seis degraus --- algumas centenas de
milhares de vezes.
\end{solution}

\begin{solution}{exo:g9:atoms-to-galaxies:6}
O \omterm{def:g9:atoms-to-galaxies:atom}{átomo} são $10^{-10}$ \unit{m} de quase nada em volta de um \omterm{def:g9:atoms-to-galaxies:atom}{núcleo}
de $10^{-15}$: empacote \omterm{def:g9:atoms-to-galaxies:atom}{átomos} numa mesa e o vazio vem junto com
eles. A solidez é a recusa elétrica das nuvens de \omterm{def:g9:atoms-to-galaxies:atom}{elétrons} a se
sobreporem --- vazio, firmemente defendido.
\end{solution}

\begin{solution}{exo:g9:atoms-to-galaxies:7}
Milímetro, $10^{-3}$; \omterm{def:g9:atoms-to-galaxies:atom}{átomo}, $10^{-10}$: sete degraus --- cerca de
dez milhões de \omterm{def:g9:atoms-to-galaxies:atom}{átomos} na fileira.
\end{solution}

\begin{solution}{exo:g9:atoms-to-galaxies:8}
\omterm{def:g8:speed-of-light:lightyear}{Ano-luz}: $10^{16}$ \unit{m} (nove mil e quinhentos milhões de
milhões). \omterm{ex:g5:stars-night-sky:milkyway}{Via Láctea}: $10^{21}$. Céu mais profundo já mapeado: perto
de $10^{26}$.
\end{solution}

\begin{solution}{exo:g9:atoms-to-galaxies:9}
$10^{5}$ ervilhas: $0.01 \times 10^{5} = \qty{1000}{m}$ --- um
``estádio'' de um quilômetro de largura: mais grandioso do que
qualquer estádio de verdade, mas a \omterm{prop:g5:mirrors-reflection:image}{imagem} se mantém --- a ervilha no
ponto central, e as ervilhas seguintes a um quilômetro de distância.
\end{solution}

\begin{solution}{exo:g9:atoms-to-galaxies:10}
Gotas nos mares: cerca de $10^{21} \times 10^{5} = 10^{26}$ --- de
modo que a velha afirmação, tomada ao pé da letra, exagerou: os mares
guardam mais gotas do que uma gota guarda \omterm{def:g7:states-of-matter:molecule}{moléculas}, por vários
degraus. A versão honesta --- as \omterm{def:g7:states-of-matter:molecule}{moléculas} de uma gota superam
qualquer coisa contável numa vida inteira --- sobrevive; as escadas
mantêm honestas até as afirmações mais queridas.
\end{solution}

\begin{solution}{exo:g9:atoms-to-galaxies:11}
Uns dez degraus para baixo até o \omterm{def:g9:atoms-to-galaxies:atom}{átomo} ($10^{0}$ a $10^{-10}$) e
vinte e seis para cima até o céu profundo --- perto do meio,
inclinado um pouquinho para o lado do pequeno. Nove anos de escola
põem quem lê ao alcance honesto das duas pontas de tudo o que já foi
medido.
\end{solution}

\begin{solution}{exo:g9:atoms-to-galaxies:12}
As respostas são pessoais --- o exercício é a própria carta. (Um
exemplar, só para dar forma: ao leitor do primeiro ano do capítulo
das \omterm{def:g1:light-and-shadows:shadow}{sombras} --- ``o gêmeo escuro na calçada um dia vai explicar os
eclipses do Sol''; ao leitor do quinto ano do capítulo da \omterm{def:g5:everyday-energy:energy}{energia} ---
``as correntes que você traça vão virar fórmulas com \omterm{def:g9:kinetic-energy-safety:joule}{joules} dentro''; ao
leitor do oitavo ano da lei de Ohm --- ``os retratos que você desenha
são o hábito de testar leis, que é o segredo inteiro''.)
\end{solution}

\begin{solution}{pb:g9:atoms-to-galaxies:1}
\textbf{1.} Pátio da escola: segundo $2$; a Terra: segundo $7$; a
\omterm{def:g6:sun-earth-moon:axisorbit}{órbita} da \omterm{def:g2:moon-in-the-sky:moon}{Lua}: segundo $9$.
\textbf{2.} Segundo $11$ (em $1.5 \times 10^{11}$ \unit{m}); o
narrador: ``a \omterm{def:g1:light-and-shadows:source}{luz} do Sol sobre a toalha tem oito minutos de idade.''
\textbf{3.} Mais ou menos do segundo $13$ ao $16$: o quadro tem o
\omterm{def:g4:solar-system:family}{Sistema Solar} inteiro como um pontinho e \omterm{def:g4:solar-system:star}{estrela} nenhuma ainda ---
honestamente, um vazio quase preto, e segundos dele: as cenas mais
verdadeiras do filme.
\textbf{4.} O céu mais profundo já mapeado, $10^{26}$ \unit{m} de
galáxias --- com a \omterm{def:g1:light-and-shadows:source}{luz} mais antiga dele mais velha que a própria
Terra.
\textbf{5.} Formiga: segundo $2$ para dentro; célula: segundo $5$;
\omterm{def:g7:states-of-matter:molecule}{molécula}: segundo $9$; \omterm{def:g9:atoms-to-galaxies:atom}{átomo}: segundo $10$.
\textbf{6.} Pelo próprio vazio do \omterm{def:g9:atoms-to-galaxies:atom}{átomo} --- cinco segundos de nada
entre a nuvem e o centro: ``se este \omterm{def:g9:atoms-to-galaxies:atom}{átomo} fosse um estádio, já teríamos
deixado as arquibancadas mais altas e ainda estaríamos voando na
direção de uma ervilha.''
\textbf{7.} O \omterm{def:g9:atoms-to-galaxies:atom}{núcleo} --- e, com honestidade: ``aqui a nossa escada
termina, não porque a natureza pare, mas porque estes são os menores
degraus que os seres humanos mediram; a construção continua.''
\textbf{8.} Para fora, $26$ segundos; para dentro, $15$: um universo
de quarenta e um segundos.
\textbf{9.} Cada segundo multiplica a vista por dez --- \emph{razões}
iguais, e não somas iguais: a lógica da escada é que um degrau é um
fator, e por fatores o passo da rua e o passo do aglomerado têm o
mesmo tamanho.
\textbf{10.} Por exemplo: ``Por mais cheio que pareça, quase tudo
aqui é vazio --- a matéria e o céu são igualmente ilhas raras num
nada imenso.''
\textbf{11.} O microscópio (capítulo das \omterm{def:g8:lenses-images:families}{lentes}: duas \omterm{def:g8:lenses-images:families}{lentes
convergentes} sobre o pequeno) e o telescópio (o mesmo capítulo,
recolhendo o fraco e distante) --- as chaves das duas direções, as
duas lapidadas na óptica deste livro.
\textbf{12.} Por exemplo: ``Nove anos, uma escada: das patas
contáveis de uma joaninha até galáxias mais velhas que o chão sob os
pés, cada degrau foi fixado por medidas que qualquer pessoa pode
conferir. No ano que vem, as perguntas que deixamos abertas --- a
\omterm{def:g2:pushes-and-pulls:force}{força} e o movimento, a natureza da \omterm{def:g1:light-and-shadows:source}{luz}, o coração do \omterm{def:g9:atoms-to-galaxies:atom}{átomo} ---
começam a ser respondidas; traga a escada.''
\end{solution}
